% !PRINTED 52
\weberpage{52}
%<<CONT
accordingly. If $\alpha_1, \alpha_2$ is a solution of this
congruence in which the two numbers are not both zero, then we obtain
all the solutions if we let the factor $h$ in $h\alpha_1, h\alpha_2$
run through the sequence of numbers $0, 1 \dots p-1$, and the
group $T$ therefore becomes, if $\alpha_1, \alpha_2$ is a fixed
pair of values satisfying the condition (9),
\[
(A_1^{\alpha_1}A_2^{\alpha_2})^h, \quad h = 0, 1, \dots p-1.
\]

The reciprocal group $U$ is obtained by seeking all the values of
the $\beta$ that satisfy the condition (9), and one thus finds
the group $U$ in the form
\[
(A_1^{\beta_1}A_2^{\beta_2})^h, \quad h = 0, 1, \dots p-1,
\]
in which $\alpha_1, \alpha_2, \beta_1, \beta_2$ are now four fixed
numbers satisfying the condition (9). If one sets $\beta_1 = 0$,
$\beta_2 = 1$, $\alpha_1 = 1$, $\alpha_2 = 0$, one obtains the
special case of the two reciprocal groups
\[
A_1^h, \quad A_2^h, \quad h = 0, 1, \dots p-1.
\]

\begin{center}
\S.\,13.

\gesperrt{The Genera in an Abelian Group.}
\end{center}

An element of an \gesperrt{Abel}ian group $S$ that is identical with its
opposite, whose second power is therefore equal to the unit
element, is called a \gesperrt{two-sided
element}\footnote{Cf. Bd.~I, S.~390, Anmerkung.}. The unit
element therefore always belongs to the two-sided ones.

If we represent the elements of $S$ by a basis
$A_1, A_2 \dots A_\nu$, whose elements have the degrees
$a_1, a_2 \dots a_\nu$, then
\weq{1}{A = A_1^{\alpha_1}A_2^{\alpha_2} \dots A_\nu^{\alpha_\nu}}
will be a two-sided element if and only if
\weq{2}{2\alpha_1 \equiv 0\ (\mathrm{mod}\ a_1), \ 2\alpha_2 \equiv 0\ (\mathrm{mod}\ a_2), \dots 2\alpha_\nu \equiv 0\ (\mathrm{mod}\ a_\nu).}
Now if among the invariants of the group a power of 2 occurs
$\lambda$ times, then we can assume the basis of $S$ so
ordered that $a_1, a_2 \dots a_\lambda$ are even and
$a_{\lambda+1} \dots a_\nu$ odd numbers. Then the solutions
of (2) come out as:
%CONT>>
